\section{Consideraciones generales}

\subsection{Sobre los servidores Windows y los registros de eventos.}
En los sistemas operativos para servidores de la familia Microsoft Windows, y de forma específica en \textbf{Windows Server 2008 R2}, la trazabilidad de la actividad del sistema se gestiona mediante la arquitectura \textit{Windows Event Log}. Dichos registros se almacenan en formato binario estructurado basado en XML con extensión \texttt{.evtx}, ubicados por defecto en el directorio protegido del sistema \texttt{C:\textbackslash Windows\textbackslash System32\textbackslash winevt\textbackslash Logs\textbackslash}. De tal manera que el resgitro de las acciones realizadas sobre el servidor se almacenan en binario y necesitamos un programa concreto para traducirlo a un texto legible por el humano.

El sistema clasifica las entradas en registros principales (\textit{Aplicación}, \textit{Seguridad}, \textit{Sistema} e \textit{Instalación}) y registros operativos específicos. Cada evento incorpora de forma estructurada un identificador numérico (\textit{Event ID}), la marca temporal exacta (almacenada internamente en formato UTC(formato estándar para el tratamiento del formato horario en los dispositivos)), el origen del evento, el nivel de gravedad, el equipo y el identificador de seguridad (\texttt{<SID>}) del usuario asociado. Desde el punto de vista pericial, estos registros constituyen una fuente primaria de auditoría que permite reconstruir cronológicamente inicios de sesión, accesos a recursos, ejecuciones de servicios y modificaciones de configuración, siempre que se verifique la integridad de los ficheros de registro y la sincronización horaria del servidor.

\subsection{Sobre las cabeceras de un correo electrónico.}
En las comunicaciones mediante correo electrónico basadas en el protocolo SMTP (estándar para la creación de mails y el estándar \href{RFC 5322}{https://www.rfc-editor.org/info/rfc5322/}(estándar del intercambio de mensajes a través de la re), todo mensaje se compone de un cuerpo y un bloque de cabeceras (\textit{headers}) que documentan su tránsito desde el origen hasta el destino. Las cabeceras contienen metadatos(grupo de datos que describen el contenido informativo de, en este caso, el correo) de enrutamiento (líneas \texttt{Received} añadidas secuencialmente en orden cronológico inverso por cada servidor de correo intermediario), direcciones de remitente y destinatario (\texttt{From}, \texttt{To}, \texttt{Return-Path}), un identificador único de mensaje (\texttt{Message-ID}), marcas temporales y resultados de validación criptográfica y de dominio siguiento los siguientes protocolos:
\begin{itemize}
	\item \texttt{SPF(\textit{Sender Policy Framework}}: lista de control de acceso publicada en los registros DNS (\textit{Domain Name System}) de un dominio donde el propietario indica qué direcciones IP (\textit{Internet Protocolo}) o servidores están autorizados para enviar correos en su nombre. En un lenguaje menos técnico, permite establecer una lista de qué personas o empresas pueden enviar un correo.
	\item  \texttt{DKIM(\textit{DomainKeys Identified Mail})}: añade una firma criptográfica digital en la cabecera del correo (\textit{DKIM-Signature}) generada con una clave privada que solo posee el servidor de salida del remitente. Añadiendo unas capas de seguridad extra, puesto que el cifrado de la comunicación emplea un mecanismo para cifrar el mensaje en tránsito y solo las persoans autorizadas pueden desencriptarlos.
	\item  \texttt{DMARC}(\textit{Domain-based Message Authentication, Reporting, and Conformance}): capa de políticas que coordina a SPF y DKIM. Exige que exista \textbf{alineación} (que el dominio visible para el usuario en el campo From: coincida con el dominio validado por SPF o DKIM) e indica al servidor receptor cómo actuar si ambos fallan.
\end{itemize}

A efectos probatorios, si bien los campos visibles en el cliente de correo (como el remitente aparente o la fecha local) son susceptibles de alteración, el análisis forense de la cadena completa de campos \texttt{Received}, las direcciones IP de los servidores intervinientes y los registros de autenticación del servidor permiten acreditar técnicamente el origen real, la ruta de transmisión y la integridad temporal del envío.

\subsection{Sobre la aplicación Dropbox y su sincronización con la nube.}
Para comprender el funcionamiento de la plataforma de almacenamiento en la nube Dropbox en entornos Windows, debe distinguirse entre el repositorio remoto en servidores del proveedor y el cliente local instalado en el equipo. El cliente sincroniza de forma automatizada una carpeta local del sistema de ficheros con la cuenta vinculada en la nube, registrando cada operación de creación, modificación, renombrado o eliminación de archivos.

A nivel forense, el cliente mantiene bases de datos locales de caché y configuración (como \texttt{filecache.dbx} o \texttt{config.dbx}) en el perfil del usuario, mientras que la plataforma remota conserva un historial de revisiones, marcas temporales de sincronización y registros de conexión (incluyendo direcciones IP y dispositivos vinculados). La correlación entre los artefactos locales del servidor y los registros de actividad de la cuenta permite establecer la trazabilidad y causalidad de las operaciones realizadas sobre los documentos sincronizados.

\subsection{Sobre el borrado de ficheros en Windows y la papelera de reciclaje.}
Para analizar la eliminación de información en \textbf{Windows Server 2008 R2}, es preciso atender al funcionamiento del sistema de archivos \textbf{NTFS} (\textit{New Technology File System}), el cual indexa cada fichero y directorio mediante registros en la Tabla Maestra de Archivos (\texttt{\$MFT}). Cuando un usuario realiza un borrado lógico en Windows, pueden producirse dos escenarios técnicamente diferenciados:

\begin{itemize}[leftmargin=1.5em]
	\item \textbf{Eliminación estándar (envío a la Papelera de Reciclaje):} El fichero no se destruye, sino que se traslada al directorio oculto del sistema \texttt{C:\textbackslash \$Recycle.Bin\textbackslash <SID>\textbackslash}, segregado por el identificador de seguridad único (\texttt{<SID>}) de cada cuenta de usuario. En dicha ubicación se generan dos artefactos: un fichero de metadatos con prefijo \texttt{\$I} (que conserva la ruta original, el tamaño y la marca temporal exacta del borrado) y un fichero con prefijo \texttt{\$R} que contiene los datos íntegros del archivo.
	\item \textbf{Eliminación permanente (\texttt{Shift + Supr}, vaciado de papelera o borrado mediante consola):} El sistema operativo marca el registro correspondiente de la \texttt{\$MFT} como disponible y libera los clústeres asignados en el mapa de bits del volumen (\texttt{\$Bitmap}), dejando el espacio como no asignado (\textit{unallocated space}) sin sobrescribir de inmediato el contenido binario en el soporte físico.
\end{itemize}

\newpage
\subsection{Sobre el borrado de ficheros en Windows y la papelera de reciclaje.}
Para analizar la eliminación de información en \textbf{Windows 11}, es preciso atender al funcionamiento del sistema de archivos \textbf{NTFS} (\textit{New Technology File System}), el cual indexa cada fichero y directorio mediante registros en la Tabla Maestra de Archivos (\texttt{\$MFT}). Cuando un usuario realiza un borrado lógico en Windows 11, pueden producirse dos escenarios técnicamente diferenciados:

\begin{itemize}[leftmargin=1.5em]
	\item \textbf{Eliminación estándar (envío a la Papelera de Reciclaje):} El fichero no se destruye, sino que se traslada al directorio oculto del sistema \texttt{C:\textbackslash \$Recycle.Bin\textbackslash <SID>\textbackslash}, segregado por el identificador de seguridad único (\texttt{<SID>}) de cada cuenta de usuario. En dicha ubicación se generan dos artefactos: un fichero de metadatos con prefijo \texttt{\$I} (que conserva la ruta original en formato Unicode, el tamaño y la marca temporal exacta del borrado) y un fichero con prefijo \texttt{\$R} que contiene los datos íntegros del archivo. Asimismo, si las carpetas del perfil se encuentran sincronizadas con \textbf{OneDrive} (integrado de forma nativa en Windows 11), la eliminación se replica en la nube trasladando el archivo a la papelera de reciclaje en línea de la cuenta.
	\item \textbf{Eliminación permanente (\texttt{Shift + Supr}, vaciado de papelera o borrado mediante consola):} El sistema operativo marca el registro correspondiente de la \texttt{\$MFT} como disponible, registra la transacción en el diario de cambios (\texttt{\$UsnJrnl}) y libera los clústeres asignados en el mapa de bits del volumen (\texttt{\$Bitmap}), dejando el espacio como no asignado (\textit{unallocated space}). En unidades de almacenamiento magnéticas (HDD) o soportes externos, el contenido binario permanece sin sobrescribir de inmediato; no obstante, en unidades de estado sólido (\textbf{SSD/NVMe}), Windows 11 emite de forma predeterminada el comando \textbf{TRIM} hacia la controladora del disco para optimizar el rendimiento, lo que puede provocar el borrado físico o desasignación inmediata de los bloques liberados.
\end{itemize}

\subsection{Sobre la recuperación de ficheros borrados en Windows}
Para proceder a la recuperación forense de archivos eliminados en sistemas \textbf{Windows 11}, debe evaluarse tanto la persistencia lógica en el sistema de archivos NTFS como la naturaleza física del soporte de almacenamiento (HDD frente a SSD con soporte \textit{TRIM}), no es lo mismo un disco magnético, como es el HDD a uno que hace la lectura por láser como el SSD y el estado del cifrado de dispositivo (\textbf{BitLocker}, habitualmente activo por defecto en equipos modernos con Windows 11) puesto que es un herramienta de Microsoft que cifra el disco en caso de detectar una actividad sospechosa sobre el mismo.

Cuando los bloques de datos no han sido sobrescritos ni purgados mediante \textit{TRIM}, la recuperación se lleva a cabo mediante técnicas de análisis de registros residuales de la \texttt{\$MFT} y del diario \texttt{\$UsnJrnl}, búsqueda por firmas de cabecera y pie de archivo en el espacio no asignado (\textit{file carving}), extracción desde instantáneas de volumen (\textit{Volume Shadow Copies} o puntos de restauración del sistema, siempre que la Protección del Sistema se encuentre habilitada en el volumen) o recuperación desde el historial de versiones y papelera remota de OneDrive. Con el fin de preservar la cadena de custodia y garantizar la validez legal de los hallazgos, toda operación de adquisición, descifrado y recuperación debe ejecutarse exclusivamente sobre una copia o imagen forense previamente verificada mediante funciones resumen criptográficas, manteniendo el soporte original inalterado.
\newpage
\subsection{Sobre los registros de log}
Todo sistema de información dispone de una serie de registros en los cuales quedan registradas todas las actividades y actuaciones que realiza el sistema, así como todos los errores que se producen al intentar realizar las diferentes tareas que tiene programadas. Estos registros se pueden almacenar en base de datos o en un fichero en el servidor.

%\subsection{Sobre la aplicación Whatsapp}
%La aplicación WhatsApp almacena sus contenidos dentro del dispositivo y estos son trasladados a la copia de seguridad que se realiza diariamente sobre el propio dispositivo.
%
%Esta copia de seguridad se almacena en la carpeta del dispositivo con nombre \texttt{Whatsapp/Databases} en diferentes ficheros con el formato ``\texttt{msgstore*.db.crypt12}''.
%
%La propia aplicación dispone de una opción llamada ``Enviar chat por correo'' que envía por correo electrónico todo el contenido de una conversación.
%
%\subsection{Sobre la conexión de dispositivos Samsung}
%Los dispositivos Samsung permiten navegar por las carpetas del dispositivo una vez que se conecta a un ordenador autorizado. Esta navegación se realiza en modo solo lectura utilizando la aplicación ``USB Write Blocker'' para evitar toda manipulación en el dispositivo.
%
%Mediante esta conexión de SOLO LECTURA se extraen las copias de seguridad almacenadas en la carpeta del dispositivo con nombre \texttt{Whatsapp/Databases} en diferentes ficheros. También mediante esta conexión de SOLO LECTURA se extraen varios ficheros de audio en formato MP3.
%
%Todos los ficheros, una vez creado su código HASH, se almacenan en un dispositivo USB con número de serie ``\texttt{xxxxxxxxxxxxxxxxxxxxx}'' que se presenta como evidencia.

\subsection{Sobre las funciones de resumen HASH.}
Los hash o funciones de resumen son algoritmos que consiguen crear a partir de una entrada (ya sea un texto, una contraseña o un archivo, por ejemplo) una salida alfanumérica de longitud normalmente fija que representa un resumen de toda la información que se le ha dado (es decir, a partir de los datos de la entrada crea una cadena que solo puede volverse a crear con esos mismos datos). Para la realización de este peritaje se ha utilizado una función HASH SHA-1.
%
%\subsection{Sobre los SMS en iPhone}
%Los SMS en iPhone se almacenan en el propio dispositivo y estos son trasladados a la copia de seguridad que se realiza con iTunes.
%
%\subsection{Sobre la herramienta iTunes}
%Se trata de una herramienta propia de Apple que sirve para realizar copias de seguridad de los dispositivos iPhone, entre otros. La copia de seguridad se realiza de manera que el dispositivo no se modifica en ningún momento.
%
%\subsection{Sobre la herramienta iPhone Backup Extractor}
%Se trata de una herramienta de la empresa Reincubate para analizar el backup de un dispositivo iPhone y extraer de forma segura su contenido sin modificar la fuente original de la información.
%
%Recupera todas tus fotos perdidas, Photostream, mensajes, SMS, videos, buzón de voz, notas, llamadas, chats de Whatsapp, contactos, entradas en el calendario y datos inaccesibles de otras aplicaciones.
%
%Funciona con las copias de seguridad de iPhone X, iPhone 8, 8 Plus, iPhone 7 Plus, iPhone 7, iPhone SE, 6 y 6S, iPhone 6 Plus y 6S Plus, iPhone 5, iPhone 5, iPhone 5c, iPhone 4, iPhone 3G, iPhone 3G, iPod, iPad de todas las versiones de iTunes, iCloud e iOS, incluido iOS 11 y 10.3.
%
%Entre los formatos de extracción que admite he extraído en PDF, HTML y CSV para los chat de Whatsapp y el VCF para los contactos.
